% 3.4 平移群的不可约表示
\section{平移群的不可约表示}

每个空间群都有一个不变子群 $\mathbf{T}$，它由空间群所依托的布拉维格子的平移对称操作构成。本节考虑如何确定布拉维格子的平移对称操作群的不可约表示。

$\mathbf{T}$ 由元素 $\{E \mid n_{1}\mathbf{t}_{1} + n_{2}\mathbf{t}_{2} + n_{3}\mathbf{t}_{3}\}$ 构成，其中 $n_{i}$ 是满足 $0 \leqslant n_{i} < N_{i}$（$i = 1, 2$ 或 3）的整数，并且施加 Born--von Kármán 循环边界条件（Born and von Kármán 1913\textit{a}, \textit{b}）。由于乘法规则形如 $\{E \mid \mathbf{t}\}\{E \mid \mathbf{s}\} = \{E \mid \mathbf{t} + \mathbf{s}\}$，又由于式 (3.1.1) 所施加的周期性，$\mathbf{T}$ 是三个阶分别为 $N_{1}$、$N_{2}$、$N_{3}$ 的循环群的外直积，于是用明显的记号可以写 $\mathbf{T} = \mathbf{T}_{1} \otimes \mathbf{T}_{2} \otimes \mathbf{T}_{3}$。循环群的表示是众所周知的。例如，$\mathbf{T}_{1}$ 的表示 $\Delta$ 共有 $N_{1}$ 个，每一个用整数 $p_{1}$（$0 \leqslant p_{1} < N_{1}$）标记，使得
\begin{equation*}
\Delta^{p_{1}}(\{E \mid n_{1}\mathbf{t}_{1}\}) = \exp(-2\pi\mathrm{i}\, n_{1}p_{1}/N_{1}).
\tag{3.4.1}
\end{equation*}
由直积群上的定理 (1.3.16)，现在可以容易地导出 $\mathbf{T}$ 的表示。它们都是一维的，共 $N_{1}N_{2}N_{3}$ 个。每个表示用三个整数 $p_{1}$、$p_{2}$、$p_{3}$（$0 \leqslant p_{i} < N_{i}$）标记。若写 $k_{i} = p_{i}/N_{i}$，则可取 $\mathbf{k} = (k_{1}, k_{2}, k_{3})$ 为倒空间中的一个矢量，即
\begin{equation*}
\mathbf{k} = k_{1}\mathbf{g}_{1} + k_{2}\mathbf{g}_{2} + k_{3}\mathbf{g}_{3}.
\tag{3.4.2}
\end{equation*}
于是 $\mathbf{k}$ 可以拿来作为 $\mathbf{T}$ 的表示的标记，这些表示就是
\begin{equation*}
\Delta^{\mathbf{k}}\left[\{E \mid n_{1}\mathbf{t}_{1} + n_{2}\mathbf{t}_{2} + n_{3}\mathbf{t}_{3}\}\right] = \exp\left\{-\mathrm{i}\,\mathbf{k} \cdot (n_{1}\mathbf{t}_{1} + n_{2}\mathbf{t}_{2} + n_{3}\mathbf{t}_{3})\right\}.
\tag{3.4.3}
\end{equation*}
现在我们可以看出，为什么在定义中把相差一个倒格子矢量的两个 $\mathbf{k}$ 矢量称为等价。因为若 $\mathbf{k}' = \mathbf{k} + \mathbf{g}$，则由于 $\exp(-\mathrm{i}\,\mathbf{g} \cdot \mathbf{t}) = 1$，
\begin{equation*}
\Delta^{\mathbf{k}'}(\mathbf{t}) = \exp\left\{-\mathrm{i}(\mathbf{k} + \mathbf{g}) \cdot \mathbf{t}\right\} = \Delta^{\mathbf{k}}(\mathbf{t}).
\tag{3.4.4}
\end{equation*}
因此，可以取任意 $N_{1}N_{2}N_{3}$ 个互不等价的 $\mathbf{k}$ 矢量作为 $\mathbf{k}$ 的允许值。上面定义的那一组就是这样一组矢量，它们的终点落在倒空间初基晶胞的内部或表面上（除原点的移动之外）；对三斜和单斜布拉维格子，我们正是用这个初基晶胞来定义布里渊区的（见定义 3.2.2）。对所有其他布拉维格子，可以取终点落在倒空间 Wigner--Seitz 晶胞内部或表面上的那一组矢量作为 $N_{1}N_{2}N_{3}$ 个非等价的 $\mathbf{k}$ 矢量（见定义 3.2.1）。这样，在任何情形下我们都可以把 $\mathbf{T}$ 的表示看作由布里渊区的点 $\mathbf{k}$ 所定义的。

有一个关于允许的 $\mathbf{k}$ 矢量密度的公式会很有用。在体积 $8\pi^{3}/V$（$V$ 是晶体晶胞的体积）中共有 $N_{1}N_{2}N_{3}$ 个 $\mathbf{k}$ 矢量。因此密度 $n(\mathbf{k})$ 为 $N_{1}N_{2}N_{3}V/8\pi^{3} = W/8\pi^{3}$，其中 $W$ 是晶体的体积。于是
\begin{equation*}
n(\mathbf{k}) = \frac{W}{8\pi^{3}}.
\tag{3.4.5}
\end{equation*}
这意味着，在无穷晶体的极限下，$\mathbf{k}$ 的取值变得任意接近，从而稠密地填满整个布里渊区。

$\Delta^{\mathbf{k}}$ 的基函数可以取 $\exp(\mathrm{i}\mathbf{k} \cdot \mathbf{r})$，因为由例 1.5.3，
\begin{equation*}
\begin{aligned}
\{E \mid \mathbf{t}\}\exp(\mathrm{i}\mathbf{k} \cdot \mathbf{r}) &= \exp\left\{\mathrm{i}\mathbf{k} \cdot (\mathbf{r} - \mathbf{t})\right\} \\
&= \exp(-\mathrm{i}\mathbf{k} \cdot \mathbf{t})\exp(\mathrm{i}\mathbf{k} \cdot \mathbf{r}) = \Delta^{\mathbf{k}}\{E \mid \mathbf{t}\}\exp(\mathrm{i}\mathbf{k} \cdot \mathbf{r}).
\end{aligned}
\tag{3.4.6}
\end{equation*}
代替 $\exp(\mathrm{i}\mathbf{k} \cdot \mathbf{r})$，我们可以使用任何形如
\begin{equation*}
\Psi_{\mathbf{k}}(\mathbf{r}) = \exp(\mathrm{i}\mathbf{k} \cdot \mathbf{r})\,u_{\mathbf{k}}(\mathbf{r})
\tag{3.4.7}
\end{equation*}
的函数，其中
\begin{equation*}
u_{\mathbf{k}}(\mathbf{r}) = u_{\mathbf{k}}(\mathbf{r} + \mathbf{t})
\tag{3.4.8}
\end{equation*}
对一切 $\mathbf{t}$ 成立。平移群 $\mathbf{T}$ 的基函数具有式 (3.4.7) 与 (3.4.8) 所定义的形式，这一结论通常称为 Bloch 定理。

\noindent\textbf{定理 3.4.1}\quad （Bloch 定理）（Bloch 1928）\ 在由 $\mathbf{T}$ 定义其周期性的周期性势场 $V(\mathbf{r})$ 中运动的粒子或准粒子的波函数可以写成
\begin{equation*}
\Psi_{\mathbf{k}}(\mathbf{r}) = \exp(\mathrm{i}\mathbf{k} \cdot \mathbf{r})\,u_{\mathbf{k}}(\mathbf{r})
\tag{3.4.7}
\end{equation*}
的形式，其中 $u_{\mathbf{k}}(\mathbf{r})$ 与 $V(\mathbf{r})$ 具有相同的周期性，而 $\mathbf{k}$ 由式 (3.4.2) 确定。\footnote{中译者注：原书在定理 3.4.1 中重复使用了编号 (3.4.7)，照录。}

然而必须强调，Bloch 函数 $\Psi_{\mathbf{k}}(\mathbf{r})$ 上的标记 $\mathbf{k}$，正如表示本身一样，只是在等价的意义下确定的，可以增减倒格矢 $\mathbf{g}$ 而改变。这是因为 $\exp(\mathrm{i}\mathbf{g} \cdot \mathbf{r})$ 满足式 (3.4.8)，所以任何这样的因子都可以吸收进 Bloch 函数的 $u_{\mathbf{k}}(\mathbf{r})$ 部分。如果所考虑的晶体置于均匀的外电场或外磁场中，理论就需要作一些修改（例如见 Ashby and Miller (1965)）。
